% Standalone chunk A.6: SPSC-Adaptive regret bound
% (POST-FIX, Round 2).
% Covers: Thm spsc_adaptive (statement) with mu = c_mu (K/T)^{1/3},
%         Rem mu_scaling (why the scaling is necessary, not cosmetic),
%         proof decomposition (false-alarm, detection delay, fixed-rate
%         allocation, exploitation+drift),
%         Rem small_change (graceful degradation for sub-threshold changes).
% Source: conf.tex lines ~477-514, 1870-1946 (post-fix).
% Shared preamble for standalone chunks. Do not compile directly; \input{} it.
\documentclass{article}
\usepackage[utf8]{inputenc}
\usepackage[T1]{fontenc}
\usepackage{lmodern}
\usepackage[margin=1in]{geometry}
\usepackage{amsmath,amssymb,amsfonts,amsthm,bm,mathtools}
\usepackage{enumitem}
\usepackage{microtype}
\usepackage{xcolor}
\usepackage{natbib}
\usepackage{booktabs}
\usepackage{hyperref}

\newcommand{\R}{\mathbb{R}}
\newcommand{\N}{\mathbb{N}}
\newcommand{\E}{\mathbb{E}}
\newcommand{\Prob}{\mathbb{P}}
\newcommand{\op}{\mathrm{op}}
\DeclareMathOperator{\sign}{sign}
\newcommand{\cA}{\mathcal{A}}
\newcommand{\cF}{\mathcal{F}}
\newcommand{\cH}{\mathcal{H}}
\newcommand{\cT}{\mathcal{T}}
\newcommand{\cI}{\mathcal{I}}
\newcommand{\cW}{\mathcal{W}}
\newcommand{\cTprobe}{\mathcal{T}_{\mathrm{probe}}}
\newcommand{\cE}{\mathcal{E}}
\newcommand{\norm}[1]{\left\lVert #1 \right\rVert}
\newcommand{\tilO}{\widetilde{\mathcal{O}}}
\newcommand{\DynReg}{\mathrm{DynReg}}
\newcommand{\rank}{\mathrm{rank}}
\newcommand{\range}{\mathrm{range}}
\DeclareMathOperator{\TopEig}{TopEig}
\DeclareMathOperator*{\argmax}{arg\,max}
\DeclareMathOperator{\tr}{tr}

\newtheorem{assumption}{Assumption}
\newtheorem{theorem}{Theorem}
\newtheorem{corollary}{Corollary}
\newtheorem{proposition}{Proposition}
\newtheorem{lemma}{Lemma}
\newtheorem{remark}{Remark}


\title{Chunk A.6: SPSC-Adaptive regret bound\\(post-fix, Round 4 verification)}
\author{(excerpt for standalone review)}
\date{}

\begin{document}
\maketitle

\textbf{Setup context.}
\begin{itemize}[leftmargin=*,itemsep=1pt,topsep=1pt]
\item SPSC-Adaptive draws probes via Bernoulli$(\mu)$ each round
(oblivious to segment boundaries); detector
$S_t=\|\widehat M_t^{\mathrm{recent}}-\widehat M_t^{\mathrm{past}}\|_\op$
compares recent-window vs.\ segment-accumulated probe moments.
\item Detector threshold calibrated so that w.p.\ $\ge 1-\delta_{\mathrm{FA}}$
no false alarm occurs under $H_0$ (no change), and all true changes
$\Delta_k=\|P_{k+1}^\star-P_k^\star\|_\op\ge 2b$ (with $b>2\eta_{\mathrm{det}}$)
are detected within a delay of $D_k=O(W_{\mathrm{det}}/\mu)$ rounds.
\end{itemize}

\textbf{What changed since Round~3 (R4 delta, this chunk).}
\begin{itemize}[leftmargin=*,itemsep=1pt,topsep=1pt]
\item \textbf{Theorem failure probability updated (R3 remaining, blocking).}
R3.1 introduced a Chernoff-based $\delta'$-event for detection delay,
but R3's theorem and ``Final.'' line still said $1-\delta-\delta_{\mathrm{FA}}$.
R4 updates both to $1-\delta-\delta_{\mathrm{FA}}-\delta'$; the theorem
now also states the hypothesis $W_{\mathrm{det}}\ge 8\log(K/\delta')$
explicitly.
\end{itemize}

\textbf{What changed since Round~2 (R3 delta, retained).}
\begin{itemize}[leftmargin=*,itemsep=1pt,topsep=1pt]
\item \textbf{Detection delay via Chernoff (R2 B3).} R2 used Wald's
identity to claim $D_k=O(W_{\mathrm{det}}/\mu)$, which gives only an
expectation bound; the theorem states a high-probability guarantee.
R3 replaces Wald with a Chernoff bound on the negative-binomial
distribution of post-change rounds, giving
$D_k\le 2W_{\mathrm{det}}/\mu$ w.p.\ $\ge 1-\delta'/K$ under
$W_{\mathrm{det}}\ge 8\log(K/\delta')$.
\item \textbf{Rem.~\ref{rem:small_change} P-drift vs.\ $\theta$-drift (R2 B4).}
R2 claimed missed sub-threshold changes are absorbed into
$O(WV)$ drift. This conflates \emph{projector drift} (not covered by
windowed ridge, since $\widehat U$ is segment-accumulated) with
\emph{parameter drift} (covered). R3 makes the distinction explicit
and offers two accounting options: (i) choose $b$ large enough that
all relevant changes are detected (giving $O(bT)$ miss cost), or
(ii) augment $V$ with the missed-change contribution inside the
$O(W\widetilde V)$ drift term.
\end{itemize}

\textbf{What changed since Round~1 (R2 delta, retained).}
\begin{itemize}[leftmargin=*,itemsep=1pt,topsep=1pt]
\item \emph{Probe-rate scaling fixed (Bug~1 fix, fatal in R1).} Old
theorem stated ``fixed $\mu\in(0,1)$'' achieves $K^{1/3}T^{2/3}$ rate;
this is \emph{false} (constant $\mu$ gives probe cost $\Theta(\mu T)$,
blowing up by $T^{1/3}$). Fix: $\mu = c_\mu\,(K/T)^{1/3}$ with a tuning
constant $c_\mu>0$.
\item \emph{Detection-delay overhead rewritten.} With the correct
scaling, $KW_{\mathrm{det}}/\mu = (c_\mu^{-1})K^{2/3}T^{1/3}W_{\mathrm{det}}$,
which is $o(K^{1/3}T^{2/3})$ iff $W_{\mathrm{det}}=o((T/K)^{1/3})$
(upper half of the R7 two-sided condition). The two-sided interval
is non-empty iff $K\log^3(K/\delta')=o(T)$; for constant $\delta'$
this is $K=o(T/\log^3 T)$ (polynomial sparsity $K=T^\beta$, $\beta<1$
is safely in; $K=T/(\log T)^c$ is in iff $c>3$).
\item \emph{Balancing argument} rewritten with Cauchy--Schwarz (not
Jensen) since under fixed rate $m_k\asymp\mu\ell_k$ (linear in $\ell_k$,
not $\ell_k^{2/3}$).
\item \emph{Rem.~\ref{rem:mu_scaling}} added to explain why the scaling
is necessary, not cosmetic.
\item \emph{Rem.~\ref{rem:small_change}} added for graceful degradation
under sub-threshold changes.
\end{itemize}

\section*{Statement}

\begin{theorem}[Adaptive SPSC]
\label{thm:spsc_adaptive}
Suppose every true subspace change satisfies $\Delta_k\ge 2b$ with the
detector separation $b>2\eta_{\mathrm{det}}$, and the detection window
satisfies the two-sided condition
$8\log(K/\delta')\le W_{\mathrm{det}} = o((T/K)^{1/3})$ for a \emph{total}
detection-delay failure budget $\delta'\in(0,1)$ (equivalently, per-change
budget $\delta'/K$ from union bound). Lower bound: Chernoff argument for
high-prob.\ delay. Upper bound: detection-delay overhead
$KW_{\mathrm{det}}/\mu$ is $o(K^{1/3}T^{2/3})$. Compatibility requires
$\log(K/\delta')=o((T/K)^{1/3})$; e.g.\ $W_{\mathrm{det}}=\Theta(\log(K/\delta'))$
saturates the lower Chernoff bound and, under compatibility, satisfies
the upper small-$o$ bound. Run SPSC-Adaptive with
Bernoulli probe rate
\[
\mu = c_\mu\,(K/T)^{1/3}
\]
for a tuning constant $c_\mu > 0$ calibrated from the same $A, B$ as
in the main Thm.\ spsc\_regret. With probability
$\ge 1-\delta-\delta_{\mathrm{FA}}-\delta'$,
\[
\DynReg_T^{(c)} = \tilO(r\sqrt{KT}) + \tilO(K^{1/3}T^{2/3})
+ O(WV) + O(T\Delta_\sigma) + O(KW_{\mathrm{det}}/\mu),
\]
where the last term is the detection-delay overhead; false alarms
occur with probability $\le\delta_{\mathrm{FA}}$ and detection-delay
excursions (some $D_k > 2W_{\mathrm{det}}/\mu$) with probability $\le\delta'$.
The adaptive variant does \emph{not} require \emph{per-segment}
lengths $\ell_k$; only the macroscopic parameters $K,T$ enter, via
both the probe rate $\mu = c_\mu(K/T)^{1/3}$ and the detection-window
admissibility $8\log(K/\delta')\le W_{\mathrm{det}}=o((T/K)^{1/3})$.
Exact segment lengths are estimated online via the change detector.
\end{theorem}

\begin{remark}[Why $\mu=\Theta((K/T)^{1/3})$ is the right scaling]
\label{rem:mu_scaling}
A fixed Bernoulli probe rate $\mu\in(0,1)$ produces $m_k\approx\mu\ell_k$
probes per segment, giving total probe cost $c\mu T$ and per-segment
subspace error $\varepsilon_k = O(1/\sqrt{\mu\ell_k})$. The
probe-plus-subspace regret is then
\[
c\mu T + \sum_k R_\cA S_w\varepsilon_k n_k
\lesssim c\mu T + \sqrt{KT/\mu},
\]
where the second term uses Cauchy--Schwarz ($\sum_k\sqrt{\ell_k}\le\sqrt{KT}$;
not Jensen, since $m_k$ scales linearly in $\ell_k$, not $\ell_k^{2/3}$).
Minimising in $\mu$: $\mu^\star = \Theta((K/T)^{1/3})$ with optimal
cost $\Theta(K^{1/3}T^{2/3})$---matching the oracle allocation of
Thm.\ spsc\_regret. A universal constant $\mu\in(0,1)$ would produce
probe cost $\Theta(T)$, blowing up the rate by $T^{1/3}$; the scaling
$\mu\asymp(K/T)^{1/3}$ is necessary, not cosmetic.
\end{remark}

\section*{Proof}

\emph{False-alarm bound ($H_0$).} Under $H_0$ (no change), triangle
inequality + matrix Freedman on each half of the detector window
gives $\Pr(S_t > 2\eta_{\mathrm{det}}) \le \delta_{\mathrm{FA}}/T$ with
$\eta_{\mathrm{det}} = C_{\mathrm{sub}}\sqrt{2\log(dT/\delta_{\mathrm{FA}})/W_{\mathrm{det}}}$
($\sqrt 2$ from each half-window having $W_{\mathrm{det}}/2$ probes).
Union over $T$ rounds: FA probability $\le\delta_{\mathrm{FA}}$.

\emph{Detection delay ($H_1$).} At $\tau_k$, separation
$\Delta_k\ge 2b > 4\eta_{\mathrm{det}}$ forces the detector once the
recent window accumulates enough post-change probes to shift the
moment by $\ge 2\eta_{\mathrm{det}}$. The number of post-change rounds
$N$ needed to collect $W_{\mathrm{det}}$ fresh Bernoulli$(\mu)$ probes is
negative-binomial with mean $W_{\mathrm{det}}/\mu$. By a Chernoff bound
for sums of Bernoulli$(\mu)$,
$\Pr(N>2W_{\mathrm{det}}/\mu)\le\exp(-W_{\mathrm{det}}/8)$. Taking
$W_{\mathrm{det}}\ge 8\log(K/\delta')$ for a \emph{total} detection-delay
failure budget $\delta'$ (equivalently, per-change budget $\delta'/K$),
the detection delay is $D_k\le 2W_{\mathrm{det}}/\mu$ w.p.\ $\ge 1-\delta'/K$;
union-bounding over $K$ changes gives total delay probability $\ge 1-\delta'$. Delay regret on this event:
$2R_\cA S_w\sum_k D_k = O(KW_{\mathrm{det}}/\mu)$.

\emph{Fixed-rate probe allocation.} With $m_k\sim\mu\ell_k$, subspace
error is $\varepsilon_k=O(1/\sqrt{\mu\ell_k})+\Delta_\sigma$.
Per-segment mismatch regret
$R_\cA S_w\varepsilon_k n_k = O(R_\cA S_w\sqrt{\ell_k/\mu})+R_\cA S_w\Delta_\sigma n_k$.
Summing with Cauchy--Schwarz:
\[
\underbrace{c\mu T}_{\text{probes}}+\underbrace{O(R_\cA S_w\sqrt{KT/\mu})}_{\text{estimation}}
+O(T\Delta_\sigma).
\]
Minimising in $\mu$: $\mu^\star=\Theta((C/c)^{2/3}(K/T)^{1/3})$ with
optimal value $\Theta(c^{1/3}C^{2/3}K^{1/3}T^{2/3})$. Calibration
$c_\mu:=(C/c)^{2/3}$ yields $\mu^\star=c_\mu(K/T)^{1/3}$ and matches
the oracle rate of Thm.\ spsc\_regret.

\emph{Detection-delay overhead under $\mu=\Theta((K/T)^{1/3})$.}
$KW_{\mathrm{det}}/\mu = (c_\mu^{-1})K^{2/3}T^{1/3}W_{\mathrm{det}}$;
this is $o(K^{1/3}T^{2/3})$ iff $W_{\mathrm{det}}=o((T/K)^{1/3})$,
which is the upper half of the theorem's two-sided condition.
Compatibility $\log(K/\delta')=o((T/K)^{1/3})$ is equivalent to
$K\log^3(K/\delta')=o(T)$; the Chernoff-admissible choice
$W_{\mathrm{det}}=\Theta(\log(K/\delta'))$ then satisfies both
bounds. For constant $\delta'$, the sharp applicability boundary is
$K=o(T/\log^3 T)$: polynomial sparsity $K=T^\beta$ ($\beta<1$) is
safely inside; $K=T/(\log T)^c$ is inside iff $c>3$ (so
$K=T/(\log T)^2$ is \emph{outside}---the theorem does not apply
there, consistent with the non-existence of an admissible
$W_{\mathrm{det}}$ in that regime).

\emph{Exploitation + drift.} Thm.\ spsc\_regret's decomposition
applies verbatim with estimated boundaries $\widehat\tau_k$; noise
contributes $\tilO(r\sqrt{KT})$, drift $O(WV)$.

\emph{Final.} Combining,
\[
\DynReg_T^{(c)}\le\tilO(r\sqrt{KT})+\tilO(K^{1/3}T^{2/3})+O(WV)+O(T\Delta_\sigma)+O(KW_{\mathrm{det}}/\mu).
\]
Total failure probability via union bound: $\delta$ (Freedman over
segment-level concentration events) $+\ \delta_{\mathrm{FA}}$ (false
alarms) $+\ \delta'$ (detection-delay Chernoff) $=\delta+\delta_{\mathrm{FA}}+\delta'$,
matching the theorem statement.\qed

\begin{remark}[Graceful degradation for small changes]
\label{rem:small_change}
If a true change at $\tau_k$ has $\Delta_k<2b$ (subspace separation
below the detector threshold), the detector may miss it, and the
algorithm's projector $\widehat U$ persists across the boundary. The
resulting per-round subspace-mismatch regret is
$R_\cA\|(I-\widehat P)\theta_t\|\le R_\cA S_w\Delta_k$, which is a
\emph{persistent linear-in-horizon} term, not a drift term. Two
accounting options: (i) if $b$ is taken large enough that all
practically relevant changes have $\Delta_k\ge 2b$ (e.g.\ $b$ matched
to the minimum actionable subspace change), missed changes
contribute negligibly, $R_\cA S_w\cdot 2b\cdot T = O(bT)$; or (ii) if
missed changes must be absorbed analytically, restate the drift term
as $O(WV+R_\cA S_w\sum_k\Delta_k\mathbf 1\{\Delta_k<2b\}\cdot\ell_k)$,
replacing $V$ with
$\widetilde V := V+\sum_k\Delta_k\mathbf 1\{\Delta_k<2b\}\cdot\ell_k/W$
inside the $O(W\widetilde V)$ drift term. In the piecewise-constant
regime this degradation is sharp: the algorithm cannot recover from a
missed subspace jump without a windowed \emph{subspace} estimator,
which SPSC-Adaptive does not deploy.
\end{remark}

\end{document}
